% Part: incompleteness
% Chapter: incompleteness-provability
% Section: fixed-point-lemma

\documentclass[../../../include/open-logic-section]{subfiles}

\begin{document}

\olfileid{inc}{inp}{fix}

\olsection{مقرر نقطے دا لمہ}

\begin{explain}
مقرر نقطے دا لمہ آکھدا اے کہ کسے وی !!{formula}~$!B(x)$
لئی اک !!a{sentence}~$!A$ اے جیہدے لئی
$\Th{T} \Proves !A \liff !B(\gn{!A})$، بشرطیکہ $\Th{T}$،
~$\Th{Q}$ دی توسیع ہووے۔ جھوٹے جملے لئی اسیں چاہواں گے
کہ $!A$ دی،~$\Th{T}$ وچ، ایس بیان نال ہم ارزی ثابت ہووے،
“$\gn{!A}$ جھوٹا اے۔” یعنی $\Gn{!A}$ اک جھوٹے
!!{sentence} دا گوڈل نمبر اے۔ ثبوت دا خیال سمجھّن لئی $!A$
دا کوائن والی غیر رسمی صورت نال ملاپ ویکھو،
“اپنا اقتباس اگے رکھن نال جھوٹ نکلدا اے” اپنا اقتباس اگے
رکھن نال جھوٹ نکلدا اے۔ کسے عبارت دے اگے اوسے دا اقتباس
رکھ کے جملہ بناؤن نوں اوس عبارت دی \emph{قطری بناوٹ}
آکھ سکدے آں۔ جیہڑا جملہ بنے اوہ اوس عبارت دی قطری صورت
اے۔ ایس طرح “اپنا اقتباس اگے رکھن نال جھوٹ نکلدا اے” دی
قطری صورت اے، “اپنا اقتباس اگے رکھن نال جھوٹ نکلدا اے”
اپنا اقتباس اگے رکھن نال جھوٹ نکلدا اے۔
دھیان دیو کہ کوائن دا جھوٹا جملہ “جھوٹ نکلدا اے” دی قطری
صورت نہیں۔ ایہ “اپنا اقتباس اگے رکھن نال جھوٹ نکلدا اے”
دی قطری صورت اے۔ سو جھوٹا جملہ بناؤن والی خاصیت وچ خود
قطری بناوٹ شامل اے!{}

حساب دی زبان وچ اک آزاد متغیر والے !!a{formula} دا اقتباس
ایویں بناندے آں۔ اوہدا گوڈل نمبر کڈھدے آں تے آزاد متغیر دی
تھاں اوس نمبر دا معیاری عددی نام رکھ دیندے آں۔
~$!E(x)$ دی قطری صورت $!E(\num{n})$ اے، جتھے
$n = \Gn{!E(x)}$۔ (اگے توں $\num{\Gn{!E(x)}}$ نوں
مختصر کر کے $\gn{!E(x)}$ لکھاں گے۔) سو جے $!B(x)$
“جھوٹا اے” والی خاصیت ہووے، تاں “اپنا اقتباس اگے رکھن نال
جھوٹ نکلدا اے” دا مطلب ہووے گا، “اپنی قطری صورت دے گوڈل
نمبر اُتے لاؤن نال جھوٹ نکلدا اے۔” فرض کرو ساڈے کول
~$\Obj{diag}$ علامت اے، فنکشن $\fn{diag}(n)$ لئی۔
ایہ فنکشن گوڈل نمبر~$n$ والے !!{formula} دی قطری صورت
دا گوڈل نمبر کڈھدا اے۔ تاں اسیں $!E(x)$ نوں
$!B(\Obj{diag}(x))$ لکھ سکدے آں۔ فیر کوائن دے جھوٹے
جملے دا روپ اوہدی قطری صورت ہووے گا، یعنی $!E(\gn{!E(x)})$
یا $!B(\Obj{diag}(\gn{!B(\Obj{diag}(x))}))$۔
~$!B(x)$ کوئی ہور خاصیت وی ہو سکدی اے تے ایہو بناوٹ
اوہدے لئی وی چلے گی۔ نامکملیت دے قضیے لئی اسیں $!B(x)$
ایہ لواں گے، “$x$،~$\Th{T}$ وچ !!{derivable} نہیں۔”
فیر $!E(x)$ دا مطلب ہووے گا، “اپنی قطری صورت دے گوڈل
نمبر اُتے لاؤن نال ایہو جیہا !!a{sentence} ملدا اے جیہڑا
~$\Th{T}$ وچ !!{derivable} نہیں۔”

ایہ گل~$\Th{T}$ وچ رسمی بناؤن لئی سانوں $\fn{diag}$ نوں
رسمی روپ دین دا طریقہ لبھنا پئے گا۔ فنکشن $\fn{diag}(n)$
قابلِ حساب اے، بلکہ بنیادی بازگشتی اے۔ جے $n$ فارمولا
~$!E(x)$ دا گوڈل نمبر اے، تاں $\fn{diag}(n)$،
~$!E(\gn{!E(x)})$ دا گوڈل نمبر دیندا اے۔ (چیتے رکھو کہ
$\gn{!E(x)}$،~$!E(x)$ دے گوڈل نمبر دا معیاری عددی نام
اے، یعنی $\num{\Gn{!E(x)}}$۔) جے $\Obj{diag}$،
$\Th{T}$ وچ فنکشن $\fn{diag}$ دی نمائندگی کرن والی علامت
ہووے، تاں اسیں $!A$ نوں فارمولا
$!B(\Obj{diag}(\gn{!B(\Obj{diag}(x))}))$ منّ سکدے آں۔
دھیان دیو کہ
\begin{align*}
\fn{diag}(\Gn{!B(\Obj{diag}(x))}) & =
\Gn{!B(\Obj{diag}(\gn{!B(\Obj{diag}(x))}))} \\
& = \Gn{!A}.
\end{align*}
فرض کرو $\Th{T}$ ایہ !!{derive} کر سکدا اے،
\[
\Obj{diag}(\gn{!B(\Obj{diag}(x))}) = \gn{!A},
\]
تاں اوہ ایہ وی !!{derive} کر سکدا اے،
$!B(\Obj{diag}(\gn{!B(\Obj{diag}(x))}))
\liff !B(\gn{!A})$۔ تعریف مطابق ایس ہم ارزی دا کھبا
پاسا~$!A$ اے۔

عام طور اُتے $\Obj{diag}$، $\Th{T}$ دی زبان دی کوئی فنکشن
علامت نہیں ہووے گی، تے~$\Th{Q}$ دی تاں یقیناً نہیں۔ پر
$\fn{diag}$ قابلِ حساب اے، ایس لئی~$\Th{Q}$ وچ اک فارمولے
$!D_{\fn{diag}}(x,y)$ نال ایہ \emph{قابلِ نمائندگی} اے۔
ایس لئی $!B(\Obj{diag}(x))$ دی تھاں
$\lexists[y][(!D_{\fn{diag}}(x,y) \land !B(y))]$ لکھ
سکدے آں۔ ایہ بدل کرن نال اُپرلا ثبوت چل پیندا اے، تے ایہ
دلیل~$\Th{Q}$ وچ ای چل پیندی اے۔
\end{explain}

\begin{lem}
\ollabel{lem:fixed-point} فرض کرو $!B(x)$ کوئی فارمولا اے
جیہدے وچ صرف متغیر~$x$ آزاد اے۔ تاں اک جملہ~$!A$ اے
جیہدے لئی $\Th{Q} \Proves
!A \liff !B(\gn{!A})$۔
\end{lem}


\begin{proof}
دتے ہوئے $!B(x)$ لئی، $!E(x)$ نوں ایہ فارمولا منّ لئیے،
$\lexists[y][(!D_{\fn{diag}}(x,y) \land !B(y))]$۔
تے $!A$~اوہدی قطری صورت اے، یعنی فارمولا $!E(\gn{!E(x)})$۔

کیوں جے $!D_{\fn{diag}}$، $\fn{diag}$ دی نمائندگی کردا اے
تے $\fn{diag}(\Gn{!E(x)}) = \Gn{!A}$، ایس لئی $\Th{Q}$
ایہ !!{derive} کر سکدا اے،
\begin{align}
  & !D_{\fn{diag}}(\gn{!E(x)}, \gn{!A}) \ollabel{repdiag1} \\
  & \lforall[y][(!D_{\fn{diag}}(\gn{!E(x)},y) \lif
  \eq[y][\gn{!A}])]. \ollabel{repdiag2}
\end{align}
ہن اسیں وکھاواں گے کہ $\Th{Q} \Proves !A \liff !B(\gn{!A})$۔
ایہ دلیل غیر رسمی طور اُتے صرف منطق تے~$\Th{Q}$ وچ
!!{derivable} حقیقتاں نال دیواں گے۔

پہلاں فرض کرو~$!A$، یعنی $!E(\gn{!E(x)})$۔ $!E(x)$ دی
تعریف مطابق $!E(\gn{!E(x)})$ ایہو فارمولا اے،
\[
\lexists[y][(!D_{\fn{diag}}(\gn{!E(x)},y) \land !B(y))].
\]
کوئی ایہو جیہا~$y$ لئیے۔ کیوں جے
$!D_{\fn{diag}}(\gn{!E(x)},y)$، ایس لئی \olref{repdiag2}
مطابق $y = \gn{!A}$۔ سو $!B(y)$ توں~$!B(\gn{!A})$ ملدا اے۔

ہن فرض کرو $!B(\gn{!A})$۔ \olref{repdiag1} نال سانوں
\begin{align*}
& !D_{\fn{diag}}(\gn{!E(x)}, \gn{!A}) \land !B(\gn{!A}).
\intertext{ایس توں ایہ نکلدا اے کہ}
& \lexists[y][(!D_{\fn{diag}}(\gn{!E(x)},y) \land !B(y))].
\end{align*}
ایہ $!E(\gn{!E(x)})$ اے، یعنی~$!A$۔
\end{proof}

\begin{digress}
ایس ثبوت دا ملاپ قابلِ حسابیت دے نظریے وچ مقرر نقطے دے لمے
دے ثبوت نال کرو۔ ایتھے اسیں اک \emph{بیان} نوں اوہدے اپنے
حوالے نال متعین کرنا چاہندے آں۔ اُتھے اسیں اک \emph{فنکشن}
نوں اوہدے اپنے حوالے نال متعین کرنا چاہندے ساں۔ ایس فرق توں
بغیر خیال اکّو اے۔
\end{digress}

\begin{prob}
  !!^a{formula}~$!A(x)$ اک \emph{سچائی دی تعریف} اے جے
  $\Th{Q} \Proves !B \liff !A(\gn{!B})$ ہر !!{sentence}s~$!B$
  لئی ہووے۔ مقرر نقطے دا لمہ ورت کے وکھاؤ کہ کوئی !!{formula}
  سچائی دی تعریف نہیں اے۔
\end{prob}

\end{document}
